\documentclass[11pt]{article}
\usepackage[T1]{fontenc}
\usepackage{lmodern,amsmath,amssymb,amsthm}
\usepackage[a4paper,margin=28mm]{geometry}
\usepackage{microtype}
\usepackage{xcolor,hyperref}
\hypersetup{colorlinks=true,linkcolor=blue!45!black,urlcolor=blue!45!black,citecolor=blue!45!black,
pdftitle={Rational-point decision problems: a conditional transfer and dependency audit},pdfauthor={Alec Kriebel}}
\newtheorem{theorem}{Theorem}
\newtheorem{lemma}[theorem]{Lemma}
\newtheorem{corollary}[theorem]{Corollary}
\newcommand{\Q}{\mathbb Q}
\newcommand{\Z}{\mathbb Z}
\newcommand{\Sha}{\operatorname{Sha}}
\newcommand{\Sel}{\operatorname{Sel}}
\newcommand{\rank}{\operatorname{rank}}
\setlength{\emergencystretch}{2em}
\title{Rational-point decision problems:\\a conditional transfer and dependency audit}
\author{Alec Kriebel\\\small ORCID: \href{https://orcid.org/0009-0001-9320-500X}{0009-0001-9320-500X}}
\date{October 6, 2026}
\begin{document}
\maketitle
\begin{abstract}
We give explicit presentation and oracle conventions for Poonen's established
transfer from rational polynomial solvability to rational-point existence on
smooth projective geometrically integral varieties over $\Q$. A hypothetical
decision procedure on this promised class would decide rational Hilbert's
tenth problem. A computable bound on the height of one rational point, when
one exists, would likewise decide existence. These implications are conditional
with respect to the proposed undecidability conclusion: the arithmetic input
in OpenAI's family 004 manuscript has not been certified by this audit.
This research document makes no unconditional solution or publication claim.
\end{abstract}

\section{Inputs and precise scope}
Let $\mathrm{H10}(\Q)$ denote the decision problem with input an integer
$n\geq0$ and a polynomial $f\in\Z[x_1,\ldots,x_n]$, asking whether
$f(\mathbf a)=0$ for some $\mathbf a\in\Q^n$. The number of variables
is part of the input. Constants, including $f=0$, are allowed.

A projective presentation $e$ consists of an integer $N\geq0$ and a finite
list of homogeneous integer polynomials $F_1,\ldots,F_r$ in
$x_0,\ldots,x_N$. It specifies
\[
 X_e=\operatorname{Proj}\bigl(\Q[x_0,\ldots,x_N]/(F_1,\ldots,F_r)\bigr).
\]
Define $\mathrm{SPGI}(\Q)$ to ask whether $X_e(\Q)\ne\varnothing$,
under the promise that $X_e$ is smooth and geometrically integral.
Projectivity is supplied by the presentation. Geometric integrality includes
nonemptiness of the geometric scheme. There is no bound on $N$, dimension,
degrees, the number of equations, or coefficient heights. The oracle need
not terminate on a presentation violating its promise.

The geometric construction and the decision transfer below are due to
Poonen~\cite{Poonen}. His Remark 1.2(c) credits a prior $\Q$-case route to
R.~Robinson, through Smory\'nski, \emph{Logical Number Theory I}, \S II.7.
We do not claim a new reduction. Family 004~\cite{OpenAI} asserts the
arithmetic undecidability input needed to turn this established implication
into an unconditional geometric conclusion. Its validity is a separate
issue; the statement appearing in a released manuscript is insufficient
certification.

\section{Poonen's effective transfer and the promise}
We use the following established input from Poonen's Theorem 1.3:
specialized to $\Q$: given a projective $\Q$-variety $X$ and an open $U\subseteq X$,
one can algorithmically construct a regular projective variety $Y$ with a
morphism $\pi:Y\to X$ whose image on rational points is exactly $U(\Q)$.
The output need not be connected or geometrically integral.

The effectivity clause matters. In Poonen's \S9, each permissible parameter
$c$ gives an output with $\pi_c(Y_c(\Q))\subseteq U(\Q)$.
An effective cardinality bound on the exceptional parameters allows a finite
list containing a parameter for which equality holds. Taking the disjoint
union of the outputs gives equality without identifying that good parameter.
The boundary exclusion comes from his Ch\^atelet-surface construction.
Resolution preserves the lifts on the smooth fibers used in the argument.
These are parts of the cited effective theorem, not a claim that naive
projective closure preserves rational points.

\begin{lemma}\label{lem:promise}
If a connected smooth projective $\Q$-scheme $C$ has a rational point,
then it is geometrically integral.
\end{lemma}
\begin{proof}
A regular Noetherian scheme has disjoint irreducible components: at an
intersection its local ring would have more than one minimal prime, whereas
a regular local ring is a domain. Its finitely many components are therefore
open and closed. Connectedness implies that $C$ is integral.
The finite algebra $L=H^0(C,\mathcal O_C)$ is consequently a field over
$\Q$. A rational point induces a $\Q$-algebra map $L\to\Q$, forcing
$L=\Q$. Global functions commute with extension of the base field here.
Over $\overline\Q$, each smooth connected projective component has
only constant functions and is integral. Thus the number of these components
is $\dim_\Q L=1$. Smoothness also gives geometric reducedness, proving
geometric integrality. This is the $\Q$-projective instance of
Poonen's Lemma 10.1.
\end{proof}

The component decomposition and geometric-integrality tests used below are
effective algebraic operations, as in Poonen's proof of Theorem 1.1(i).
For clarity, for a smooth projective $Y$, effective coherent cohomology
computes the finite \emph{\'etale} algebra $H^0(Y,\mathcal O_Y)$.
Its field factors give the connected components, and a connected component
is geometrically integral precisely when its global-functions field has
$\Q$-dimension one. Projective constructions provide embeddings, and
elimination gives homogeneous-equation presentations. Regularity and
smoothness agree for finite-type schemes over the perfect field $\Q$.
These conventions ensure that no invalid input is passed to the oracle.

\begin{theorem}\label{thm:transfer}
An algorithm deciding $\mathrm{SPGI}(\Q)$ on every valid input would
decide $\mathrm{H10}(\Q)$. Conversely, decidability of
$\mathrm{H10}(\Q)$ would decide $\mathrm{SPGI}(\Q)$.
\end{theorem}
\begin{proof}
For the forward implication let $D$ be the promised decision procedure.
On input $f$, handle constant polynomials directly. Otherwise form the
reduced affine zero scheme $U$ of $f$ and its projective closure $X$,
recording $U$ as an open in $X$. If $U$ is empty, return NO.
Apply Poonen's effective theorem to obtain $(Y,\pi)$.
Compute the connected components $C_1,\ldots,C_s$ of $Y$ and their
projective presentations. Test geometric integrality and discard each
component that fails. By Lemma~\ref{lem:promise}, discarded components
have no rational points. Every retained component is smooth, projective
and geometrically integral. Query $D$ on each retained presentation,
and return YES if and only if some query answers YES.

Every preprocessing step terminates. The list of valid queries is finite,
so $D$'s promise ensures termination of the whole algorithm. Correctness
follows from
\[
 f(\Q^n)\ni0
 \ \Longleftrightarrow\ U(\Q)\ne\varnothing
 \ \Longleftrightarrow\ Y(\Q)\ne\varnothing
 \ \Longleftrightarrow\ \exists i:\ C_i(\Q)\ne\varnothing.
\]
The last condition is exactly the disjunction of the retained oracle answers.
If there are no retained components it is false.

Conversely, on a projective presentation take the $N+1$ standard affine
charts, setting $x_j=1$ on chart $j$. Their equations have a common rational
zero precisely when the sum of their squares does, since $\Q$ is ordered.
Call an $\mathrm{H10}(\Q)$ decision procedure on these finitely many
polynomials and take the disjunction. Every rational projective point lies
in one of these charts, proving the reverse implication.
\end{proof}

The forward construction yields a finite nonadaptive disjunctive oracle
reduction, and in particular a Turing reduction. It does not establish a
single-output many-one reduction into the promised geometrically integral
class. No quantitative or dimension-preservation assertion is used.

\begin{corollary}\label{cor:conditional}
If $\mathrm{H10}(\Q)$ is undecidable, then no algorithm decides
rational-point existence for the presented class $\mathrm{SPGI}(\Q)$.
\end{corollary}
This corollary is conditional in the present audit. Its arithmetic antecedent
must be verified independently before claiming that a hypothesis has been removed.

\section{An explicit conditional height consequence}
Specify a size parameter as follows. For $a\geq0$, let $b(a)$ be its binary
representation, with $b(0)=0$, and let $\ell(a)$ be its length. Use the
self-delimiting integer code $c(a)=1^{\ell(a)}0b(a)$.
Encode $N$, $r$, then the number of monomials in each polynomial. For every
monomial encode a sign bit, $c$ of the nonzero coefficient magnitude,
and $c$ of each of its $N+1$ exponents. Require increasing lexicographic
exponent vectors and a common total degree within each polynomial.
The zero polynomial has zero monomials. These conventions specify a
uniquely parsable finite string $e$; let $s(e)$ be its bit length.
The size refers to this presentation and embedding, not to the isomorphism
class of the abstract variety.

For $P\in\mathbf P^N(\Q)$ choose its primitive integer coordinates
$(z_0,\ldots,z_N)$ with the first nonzero coordinate positive, and set
$H(P)=\max_j|z_j|$.

\begin{corollary}
Assume $\mathrm{H10}(\Q)$ is undecidable. There is no total computable
function $B:\mathbb N\to\mathbb N$ such that for every valid $e$ with
$X_e(\Q)\ne\varnothing$ there exists $P\in X_e(\Q)$ satisfying
$H(P)\leq B(s(e))$.
\end{corollary}
\begin{proof}
Compute $M=B(s(e))$. Enumerate all nonzero primitive integer vectors in
$[-M,M]^{N+1}$ with first nonzero coordinate positive, and evaluate the
$F_i$ exactly. Return YES if any vector satisfies every equation, and NO
otherwise. The enumeration is finite. The hypothesized bound puts a point
in the enumeration whenever one exists. This would decide
$\mathrm{SPGI}(\Q)$, contradicting Corollary~\ref{cor:conditional}.
For $M=0$ the enumeration is empty, consistently with $H(P)\geq1$.
\end{proof}

The proof also rules out a presentation-dependent computable bound that
terminates on every valid input. It concerns the height of \emph{one} point
when points exist. It makes no assertion about a fixed-dimensional family,
finite rational-point sets, or all points on one variety.

\section{Arithmetic validation gap and boundaries}
Family 004's asserted main theorem has exactly the variable-count input
convention above. Its proof uses finite rational tests, compactness,
additive elliptic indices and arithmetic height estimates. The current
audit has reconstructed the conditional transfer and several arithmetic
interfaces, but has not certified all pivotal new companion arguments.
No applicable complete Lean formalization of family 004 has been identified;
this is a verification limit, not evidence that the theorem is false.

One particularly precise remaining dependency occurs in its pole-parity
argument. For the special curves
\[
 E_l:\quad y^2=x(x-l)(x+3l),
\]
with $l$ restricted to the parameters selected by family 004's descent
construction, that construction gives $\dim_{\mathbf F_2}\Sel_2(E_l/\Q)=3$ and
full rational two-torsion. Hence
\[
 r+\dim_{\mathbf F_2}\Sha(E_l/\Q)[2]=1,
 \qquad r=\rank E_l(\Q).
\]
The desired realization of its extra Selmer class by a rational point
follows if $\Sha(E_l/\Q)[2^\infty]$ is finite: the alternating perfect
Cassels--Tate pairing then forces even dimension of $\Sha[2]$, yielding
$r=1$ and $\Sha[2]=0$. Low Selmer corank alone does not supply that
finiteness. The cited new pointwise $2$-converse~\cite{TwoConverse}
is the upstream assertion intended to do so and remains under audit.
This document does not assume its conclusion by silently replacing it with
a named conjecture or a mismatched formalization.

Naive compactification cannot replace the effective transfer:
$\operatorname{Spec}\Q[x,y]/(x^2+1)$ has no rational point, whereas
its closure $V(X^2+Z^2)\subseteq\mathbf P^2$ contains $[0:1:0]$.
Resolution alone need not preserve point existence in the other direction:
$V(x^2+y^2)$ has a rational singular point but its normalization is
$\mathbf A^1_{\Q(i)}$, which has no rational point. A morphism from a
resolution cannot create a rational point when its target has none.
The issue is losing singular points or retaining boundary points introduced
by compactification.

No conclusion for curves, surfaces, fixed dimension, Fano varieties,
general type, or a finite-point promise follows automatically. This audit
asserts neither a Turing degree nor quantitative bounds beyond the stated
oracle implication. The height corollary is an elementary computability
consequence and adds no independent arithmetic input.

\paragraph{Disclosure and status.}
AI tools were used extensively in research, drafting and adversarial
verification. Those automated reviews are not conventional human peer review.
This is a local conditional research document; it has not been deposited
as an unconditional solution or conventionally refereed.

\begin{thebibliography}{9}
\raggedright
\bibitem{Poonen} Bjorn Poonen,
\emph{Existence of rational points on smooth projective varieties},
J. Eur. Math. Soc. \textbf{11} (2009), 529--543.
\href{https://doi.org/10.4171/JEMS/159}{doi:10.4171/JEMS/159}.
\bibitem{OpenAI} OpenAI,
\emph{Hilbert's tenth problem over the rational numbers},
OpenAI Math Release preprint, September 24, 2026.
Manuscript key: OAI:Hilberts-tenth-problem-over-the-rational-numbers-September-24-2026.
Source version: commit \texttt{adc7f1241b42e322a6451854ab7e4b4c146bf78a}.
\href{https://github.com/openai/math/tree/adc7f1241b42e322a6451854ab7e4b4c146bf78a/preprints/Hilberts-tenth-problem-over-the-rational-numbers-September-24-2026}{Pinned source and citation metadata}.
\bibitem{TwoConverse} OpenAI,
\emph{A pointwise $2$-converse for elliptic curves with rational two-torsion},
OpenAI Math Release preprint, September 24, 2026, same pinned commit.
\href{https://github.com/openai/math/tree/adc7f1241b42e322a6451854ab7e4b4c146bf78a/preprints/A-pointwise-2-converse-for-elliptic-curves-with-rational-two-torsion-September-24-2026}{Pinned source and citation metadata}.
\end{thebibliography}
\end{document}
